\documentclass[namedate,webpdf,imaiai]{ima-authoring-template}

\usepackage{booktabs}
\usepackage[nameinlink,noabbrev]{cleveref}

\providecommand{\PaperCommented}{0}
\newcommand{\draftnote}[1]{%
  \ifnum\PaperCommented=1
    \par\medskip\noindent\textbf{TODO:} #1\par\medskip
  \fi
}

\title[Beyond first attendance]{Beyond First Attendance: Interpreting Visit Frequency in University Mathematics Support}
% Current author block retained for restoration at submission stage.
% \author{Heriberto Espino-Montelongo*\ORCID{0009-0009-1230-2931}
% \address{Universidad de las Am\'ericas Puebla}}
% \author{Daniela Cort\'es-Toto
% \address{Universidad de las Am\'ericas Puebla}}
% \authormark{Espino-Montelongo and Cort\'es-Toto}
% \corresp[*]{Corresponding author: \href{mailto:heriberto.espinomo@udlap.mx}{heriberto.espinomo@udlap.mx}}

\author{Authors to be confirmed\address{Affiliations pending institutional review}}
\authormark{Authors pending}

\DOI{}
\copyrightyear{2026}
\vol{00}
\pubyear{2026}
\firstpage{1}

\abstract{
Recorded attendance at university mathematics-support centres may reflect academic help-seeking and incentives arising from other university programmes, complicating its interpretation as a dose of instruction. At Universidad de las Am\'ericas Puebla, first-semester students can complete one activity in a social and cultural integration programme (PPA1) by making three visits to the mathematics support centre (CMAT); this option was available throughout the 11 study periods. We examine 6,627 first eligible university mathematics-course attempts across 190 instructor--period contexts, with outcomes of standardised final performance ($Z$) and passing (PASS, numeric grade $\geq7.5$). Models adjust for instructor--period and degree programme; pairwise Wald comparisons use Holm multiplicity correction. Any recorded attendance was associated with 0.361 standard deviations higher adjusted $Z$ (95\% CI 0.302--0.421) and a 15.3-percentage-point higher adjusted PASS probability (95\% CI 12.4--18.1). In the outcome-blind support-based user-only grouping, omnibus tests were not significant ($p=0.317$ for $Z$, $p=0.152$ for PASS), and no positive-frequency pair survived Holm adjustment. In a finer exploratory model including non-users, $7+$ versus one visit was associated with a 14.5-percentage-point higher PASS probability (Holm-adjusted $p=0.0206$), but the corresponding user-only comparison was not Holm-significant. The ranges $0$, $1$--$3$, $4$--$6$ and $7+$ are descriptive, not four confirmed outcome regimes. Our contribution is to distinguish initial attendance from the limited inferential resolution of repeated use, in a setting where a continuously available integration incentive may affect recorded visits. The analyses do not identify causal effects of either attendance or PPA1.
}

\keywords{mathematics support centre; mathematics learning support; visit frequency; Holm correction; multiple comparisons; academic performance; higher education}

\begin{document}
\maketitle
\begin{center}\small\emph{Provisional manuscript; authorship, ethics, PPA empirical tests and administrator scientific approval remain pending.}\end{center}

\section{Introduction}

University mathematics and statistics support centres provide supplementary learning opportunities for students entering programmes with different levels of prior mathematical preparation. Reviews have documented the development of such services and the methodological difficulty of evaluating them with attendance records alone \citep{matthews2013,lawson2020,mullen2024}. Participation is a decision made by students within an academic and institutional context rather than an assigned intervention. A centre can be used to resolve an immediate question, repeatedly revisit difficult material, or satisfy a participation activity unrelated to mathematics coursework; consequently, equal visit counts need not represent equal instructional need or intensity. At this institution PPA1 is a first-semester social and cultural integration programme, in which activities can include attending sports events, theatre or dance performances or participating in student organisations. Completing three CMAT visits is one available PPA1 activity, offered throughout the study period, although actual credit receipt and individual eligibility are not observed for the complete MU cohort.

Prior research on university mathematics-support centres demonstrates that the difference between initial attendance and repeated participation has been an empirical concern for more than a decade. At the National University of Ireland Maynooth, \citet{macanbhaird2009} compared first-year students who attended the mathematics support centre more than once with those who attended once or never; among Arts and Finance students, the former group had an 89\% pass rate compared with 62\% for the latter, whereas the Science-student comparison was considerably weaker. Their use of Leaving Certificate mathematics achievement and diagnostic assessment also highlighted the importance of prior mathematical preparation, but their combined zero-or-one comparison does not resolve whether a first visit differs from non-attendance.

A more detailed distinction was made by \citet{jacob2019} using twelve years of administrative data from Dublin City University. Their logistic model controlled for previous mathematical achievement and module studied and reported passing odds ratios, relative to never attending, of 1.63 for a single visit and 13.78 for fifteen or more visits. These are associations in the \emph{odds} of passing, rather than percentage-point effects or causal returns to additional visits. In a US Calculus I tutoring-centre study, \citet{rickard2018} used regression adjustment for prior ability and found that three tutoring visits were associated with approximately one percentage-point higher course grades. Importantly, results are not uniformly favourable: a smaller US drop-in tutoring pilot by \citet{navarramadsen2010} found similar grades for mathematics majors who used the service and those who did not, with weak correlations between tutoring hours and grades. Together, these studies motivate distinguishing the service-use decision from the frequency profile rather than assuming a common graded benefit across institutions.

Evidence from Latin America provides an additional relevant comparison. At the Instituto Tecnol\'ogico de Costa Rica, \citet{chaconvargas2022} analysed participation in a first-year General Mathematics peer-tutoring programme with 528 enrolled participants in 2019. Although mean grades tended to be higher among students with more sessions, tests of grade distributions generally found no statistically significant differences except for a group attending at least fourteen of sixteen sessions; only 26 students belonged to that high-frequency group. This setting involves scheduled peer-tutoring and an additional preparatory programme rather than CMAT's drop-in faculty support, but it illustrates the interpretive challenges of frequency categories and small upper-tail groups in a closer regional context.

Differences between support centres are a further reason for caution when transporting observed associations. \citet{byerley2024} assembled information from over 26,000 students at ten US mathematics tutoring centres, combining records of visits, prior preparation and grades with organizational characteristics to formulate hypotheses about centre structure. Their study did not formally establish the effectiveness of particular organizational features. The broader evaluation review by \citet{mullen2024}, which included 148 publications from twelve countries, likewise documents substantial differences in support formats, data sources and evaluation designs. \citet{matthews2013} and \citet{lawson2020} provide methodological and institutional perspectives on the development and evaluation of such services.

Attendance is behavioural rather than assigned. Surveys of students who do not engage with mathematics support document heterogeneous practical, academic and informational reasons for non-use \citep{macanbhaird2013}. A postsecondary help-seeking meta-analysis finds that the association with achievement depends on the type of help sought \citep{fong2023}, whereas a centre-entry count cannot distinguish those forms of engagement. Experimental encouragement studies also clarify the inferential boundary. Randomised messages increased take-up of peer tutoring without establishing a corresponding improvement in grades \citep{pugatch2018}; a randomised encouragement design for supplemental instruction similarly produced a clear first-stage participation response but imprecise academic-performance estimates \citep{paloyo2016}. These interventions concern encouragement, not random assignment to the amount of mathematics tutoring received.

There is also a distinction between studying \emph{why students attend} and studying their \emph{final outcomes conditional on attendance}. Using drop-in provision at Coventry University, \citet{gokhool2026} separately models whether students engage at all and the number of visits among those who engage; that study is about predictors of service use, not the effect of each visit on grades. Our data ask the complementary academic-outcome question. For positive attendance frequencies, testing every possible pair becomes difficult when the upper tail contains relatively few students and nominal $p$-values invite selective interpretation. The analysis therefore treats initial engagement and frequency among users as separate inferential objects, preserves an outcome-blind grouping decision and adjusts the full declared pairwise family using Holm's procedure.

This paper examines the association between recorded CMAT visit frequency and final outcomes in the first eligible attempt at a university mathematics course at Universidad de las Am\'ericas Puebla, Mexico. It separates the benchmark comparison of no recorded attendance versus any attendance from the central analysis among students who used the centre. We examine two outcomes: a within-instructor--period-standardised final-grade measure and the probability of passing. For an eligible first-semester student, a third CMAT visit can complete one PPA1 activity, while subsequent attendance goes beyond the requirement of that particular option. The visit count does not reveal whether any individual attended to obtain PPA1 credit. The primary inference uses instructor--period and degree-programme adjustment, cluster-robust omnibus tests and pairwise Wald contrasts whose multiplicity is controlled by the Holm procedure. The positive-frequency grouping originally selected from an outcome-blind support audit is reported separately from a more finely resolved exploratory attendance grid.

Our research questions ask whether adjusted outcomes differ between no use and any recorded use; whether positive frequencies can be distinguished after multiplicity correction; and how far the descriptive bands of zero, one to three, four to six and seven or more visits can be interpreted without claiming four empirically homogeneous groups. The $3|4$ boundary marks the first visit beyond a possible PPA1 activity for eligible students, not a proven academic threshold; $7+$ is an exploratory tail grouping. We assess sensitivity to grouping and instructor-level clustering. The study cannot identify a minimum effective visit count, measure the causal impact of PPA1 incentives or establish causal returns to attendance.
